% Pista ana-25i-q4a · Anàlisi
% Procedència: PAU Catalunya 2025 · Convocatòria de juny (incidències) · Sèrie 4 · Exercici 4 — Opció A
\documentclass[11pt,a4paper]{article}
\usepackage{pista}
\pistainfo{Catalunya 2025}{Convocatòria de juny (incidències)}{Sèrie 4}

\begin{document}

\section*{\textcolor{blauPista}{Exercici 4 --- Opció A}}

\noindent El dissenyador d'una casa de perfums ha decidit que l'ampolla del proper perfum que l'empresa traurà al mercat serà de vidre i tindrà la forma d'un cilindre amb una mitja esfera a la part superior, tal com es pot veure a la figura següent:

\noindent\textit{(A la figura: $r$ és el radi del cilindre i de la mitja esfera, i $h$ és l'altura del cilindre.)}

\noindent El perfum estarà contingut a la part cilíndrica, la qual té una capacitat de $100\,\text{cm}^{3}$. L'enginyer a càrrec del procés de producció es proposa minimitzar el cost total del vidre usat i, per tant, vol que la superfície total de l'ampolla sigui mínima.

\subsection*{\textit{a)} \normalfont Comproveu que la superfície total de l'ampolla de perfum ve donada per la fórmula
\[
S(r) = 3\pi r^{2} + \dfrac{200}{r}.
\]
\hspace*{\fill}\textnormal{[1~punt]}}

\pista{En primer lloc, el volum d'un cilindre es calcula seguint aquesta fórmula:
\[
V = \pi r^{2} \cdot h,
\]
on $r$ és el radi de la base i $h$ és l'altura. Imposa que el volum valgui $100$ i aïlla $h$.

\medskip
A continuació, la superfície de l'ampolla té tres parts: la base del cilindre, que és un cercle; la cara lateral del cilindre; i la mitja esfera de dalt. Fixa't que la tapa de dalt del cilindre no hi és, perquè la cobreix la mitja esfera. La cara lateral, si la desplegues, és un rectangle: la base és la longitud de la circumferència, $2\pi r$, i l'altura és $h$. Per a la mitja esfera, fes servir la fórmula de la \textit{Nota} de l'apartat \textit{b)}, $4\pi r^{2}$, i agafa'n la meitat.

\medskip
Finalment, suma les tres parts i substitueix-hi $h$. Arribaràs a l'expressió $S(r)$ de l'enunciat.}

\subsection*{\textit{b)} \normalfont Quines dimensions ha de tenir l'ampolla de perfum perquè la superfície total sigui mínima? Quina és la superfície total de l'ampolla en aquest cas? \hspace*{\fill}\textnormal{[1,5~punts]}}

\noindent\textit{Nota:} Recordeu que la superfície d'una esfera ve donada per la fórmula $4\pi r^{2}$, on $r$ és el seu radi.

\pista{En primer lloc, deriva $S(r)$. Recorda que si escrius $\frac{200}{r}$ en forma de potència, $200\,r^{-1}$, pots derivar-la fàcilment. A continuació, resol $S'(r) = 0$. Per treure la $r$ del denominador, multiplica els dos costats de l'equació per $r^{2}$. Quan arribis a $r^{3}$ igual a un nombre, fes-ne l'arrel cúbica amb la calculadora. Comprova que realment és un mínim relatiu (estudia el signe d'$S''$ en aquest punt).

\medskip
Finalment, tingues en compte que la teva resposta ha d'incloure, d'una banda, les dues mides de l'ampolla, el radi $r$ i l'altura $h$; i d'altra banda, la superfície total. Per a l'altura $h$, fes servir l'expressió que has trobat a l'apartat \textit{a)}.}

\end{document}
